\begin{appendix}

\section{Proofs}\label{app:proofs}

\begin{proof}[Proof of Theorem~\ref{thm:main}]
All probability bounds for the endpoints are uniform over \(a\in\mathcal C_a\).
Write \(\mathcal D_n\) for the training data and the fitted nuisance functions,
including their training randomization. Uniform boundedness, the fixed
two-dimensional rectangle class, and relative VC Bernstein bounds give
\[
r_n=\sqrt{\frac{\log(B_ns_n)}{k_n}}+
\frac{\log(B_ns_n)}{k_n}+
\|\widehat m-m_a\|_\infty=o_p(1).
\]
These bounds hold simultaneously for the root splitting and estimation
samples and all their rectangular restrictions. They therefore cover both
ordinary and failure trees, even though the latter use conditional transition
laws. Increasing recodings do not enlarge this rectangle class. The score
identity and its perturbation bound are
\[
\E(\Gamma_i\mid X_i=x)=\tau_a(x),
\qquad
\max_i|\widehat\Gamma_i-\Gamma_i|
\le2\|\widehat m-m_a\|_\infty.
\]
Combining the concentration bounds with \eqref{eq:transfers} yields
\begin{equation}\label{eq:proof-criterion}
\sup_{C,j,t}
\left|\widehat{\mathcal Q}^{L}_C(j,t)-\mathcal Q_C(j,t)\right|
=o_p(1),
\end{equation}
where \(\mathcal Q_C\) is the population score gain.

First consider a node before coordinate \(j\) has been used. Its range in
that coordinate is still \([0,1]\), and independence removes the other
component from the split contrast. Direct calculation gives
\begin{equation}\label{eq:proof-shape}
\mathcal Q_C(j,t)=
\begin{cases}
a_j^2t/(1-t),&0<t\le1/2,\\
a_j^2(1-t)/t,&1/2\le t<1.
\end{cases}
\end{equation}
The unique maximizer is \(1/2\), with gain \(a_j^2\). Positivity of
\(\underline a\) and the coefficient gap imply
\(a_2^2-a_1^2\ge2\underline a\Delta_a>0\).
At nodes with at least \(C_0k_n\) splitting observations, cuts approaching
\(1/2\) are size-admissible and balanced on the concentration event.
Maximization and \eqref{eq:proof-criterion} thus imply
\begin{equation}\label{eq:proof-localization}
\max_{j,F}|\widehat t_{j,F}-1/2|=o_p(1),
\end{equation}
where \(F\) ranges over first-use nodes above this size band in the ordinary
and failure trees.

The continuation condition now supplies the depth bound. It is used here
in addition to balance. If \(n_d(x)\) is the splitting-sample size on a path,
then, whenever a split is made,
\[
\beta n_d(x)\le n_{d+1}(x)\le(1-\beta)n_d(x).
\]
The root size is proportional to \(s_n\), no path stops above \(C_0k_n\),
and every child has at least \(k_n\) observations. Consequently there are
constants \(0<c_\beta<C_\beta<\infty\) and an integer \(K\), independent of
\(n,a,h\), such that, with probability tending to one,
\begin{equation}\label{eq:proof-depth}
c_\beta H_n\le J_{h,n}(X),J_{h,n}^0(X)\le C_\beta H_n.
\end{equation}
Here \(J_{h,n}\) is the number of candidate draws on an ordinary path,
and \(J_{h,n}^0\) is defined when \(\omega_n<1\).
At most \(K\) draws occur after the path enters the band
\(n_d<C_0k_n\), including a final unsuccessful search. Reached leaves have
splitting sizes between \(k_n\) and \(C_0k_n\). Relative concentration gives
their population masses of order \(k_n/s_n\) and honest estimation counts
of order \(k_n\). In particular,
\begin{equation}\label{eq:proof-leaf}
\max_{b,L}\left|
\frac{1}{N_{b,n}^{\rm est}(L)}
\sum_{i\in I_{b,n}^{\rm est}:X_i\in L}\widehat\Gamma_i
-\E\{\tau_a(X)\mid X\in L\}
\right|=o_p(1).
\end{equation}
The maximum includes both representations. The counts are positive on this
event.

We next control the approximation after first use. Choose a deterministic
\(\delta_n\downarrow0\) so that \eqref{eq:proof-localization} is at most
\(\delta_n\) with probability tending to one. For each coordinate, its
first-use nodes above the terminal band form an antichain. Label only their
children by the sign suggested by the side of the cut, and retain this label
in their descendants. The total population mass receiving a wrong label is
\[
\sum_F P_X(F)|\widehat t_{j,F}-1/2|\le\delta_n,
\]
where the sum is over these above-band first-use nodes.
For a descendant \(C\), let \(q(C)\) be the conditional fraction with the
wrong label. Then
\[
\sup_t\mathcal Q_C(j,t)\le4\overline a^2q(C),
\qquad
\int_C|\E(s_j\mid C)-s_j(x)|\,dP_X(x)\le4P_X(C)q(C).
\]
These bounds also apply after further splits. For
\(u_n=\sqrt{\delta_n}\), select the first descendant on each path for which
\(q(C)>u_n\). Those descendants are disjoint, so their combined mass is at
most \(\delta_n/u_n\). Outside this vanishing set, a previously used
coordinate has gain at most \(4\overline a^2u_n\) at every subsequent node.
An offered unused coordinate has gain at least
\(\underline a^2-o_p(1)\) above the terminal band and therefore wins.
Only the root can involve competition between two unused coordinates.

The repeated coordinate is offered with probability at least \(1/2\) at
each draw. Together with \eqref{eq:proof-depth}, the preceding ordering
argument bounds its probability of remaining unused above the terminal
band by
\[
2^{-\lfloor c_\beta H_n\rfloor+K+1}+o(1)=o(1).
\]
For the singleton, the probability of first use only in that band is also
\(o(1)\). If \(\omega_n\le H_n^{-1/2}\), it is at most \(K\omega_n\).
Otherwise, avoiding an earlier use has probability at most
\(\exp\{-\omega_n(\lfloor c_\beta H_n\rfloor-K-1)\}+o(1)\).
The error terms here include the exceptional target mass just bounded.

Let \(D_{h,n}(x)\) indicate first use of the singleton in a generic tree,
and let \(\check T^{\rm tree}_{h,n}\) be that tree's honest score prediction.
If \(D_{h,n}(x)=0\), its leaf has no restriction in the singleton coordinate,
whose conditional mean is exactly zero. For a first use above the terminal
band, localization and the wrong-label bound control the integrated
approximation error. A late first use adds at most twice its probability to
the integrated error for the sign component, which vanishes by the preceding
bound.
The same bound applies to the repeated component. Thus
\begin{equation}\label{eq:proof-mixture}
\E_\Theta\!\left[
\left\|\check T^{\rm tree}_{h,n}
-D_{h,n}\tau_a-(1-D_{h,n})a_hs_h\right\|_{L^1(P_X)}
\middle|\mathcal D_n\right]=o_p(1).
\end{equation}
This is an approximation after first use, not an assertion that every leaf
is exactly pure.

Suppose now that \(\omega_n<1\). Conditional on \(Q_{M_n}=q\), the singleton
is included with probability \(q/M_n\), so its unconditional inclusion
probability at a node is \(\omega_n\). On the locally conditioned failure
tree, write \(\rho_{h,n,t}^0(x)\) for the probability that the ordinary
transition at the current node chooses the singleton as its splitting
coordinate. It satisfies \(0\le\rho_{h,n,t}^0\le\omega_n\).
After the root and above the terminal band, this probability equals
\(\omega_n\) outside the exceptional paths. With
\(\Lambda_{h,n}^0=\sum_{t\le J_{h,n}^0}\rho_{h,n,t}^0\), we obtain
\begin{equation}\label{eq:proof-hazard}
\left|\Lambda_{h,n}^0-J_{h,n}^0\omega_n\right|
\le(K+1)\omega_n
\end{equation}
outside a set of joint training, failure-tree, and test-point probability
tending to zero.

Node locality ensures that conditioning transitions on other branches leaves
the transition kernels on the target path unchanged. At each stage, factor
the probability of avoiding a singleton split from
the transition kernel conditional on that event. Multiplication along the
path and integration over the conditional kernels give the exact identity
\begin{equation}\label{eq:proof-survival}
u_{h,n}(x):=\Prb_\Theta\{D_{h,n}(x)=0\mid\mathcal D_n\}
=\E_{\Theta^0}\!\left[
\prod_{t=1}^{J_{h,n}^0(x)}(1-\rho_{h,n,t}^0(x))
\middle|\mathcal D_n\right].
\end{equation}
The kernels are well-defined because their conditioning probabilities are
at least \(1-\omega_n>0\). This identity concerns no first use throughout
the path, exactly the event defining \(D_{h,n}=0\).

If \(\mathcal I_n\to0\), then \(\omega_n\to0\) and
\(\Lambda_{h,n}^0\to_p0\) by \eqref{eq:proof-depth} and
\eqref{eq:proof-hazard}. The product converges to one since its deficit is
at most \(\Lambda_{h,n}^0\). If \(\mathcal I_n\to\infty\) and
\(\omega_n<1\), the cumulative hazard diverges and
\[
\prod_t(1-\rho_{h,n,t}^0)\le e^{-\Lambda_{h,n}^0}\longrightarrow_p0.
\]
For indices with \(\omega_n=1\), both coordinates are always offered.
The stronger one wins at the root and the other at the next interior node,
outside the same exceptional set. Thus recovery also holds on these indices,
without a failure-tree conditioning on a null event.

For the intermediate case, fix \(a\in\mathcal C_a\).
The condition \(\mathcal I_n\to\kappa\in(0,\infty)\) implies
\(\omega_n\to0\). Equations~\eqref{eq:intermediate-depth} and
\eqref{eq:proof-hazard} yield
\(\Lambda_{h,n}^0\to_p\kappa d_h(a)\). Moreover,
\[
0\le-\log\prod_t(1-\rho_{h,n,t}^0)-\Lambda_{h,n}^0
\le\frac{\omega_n\Lambda_{h,n}^0}{1-\omega_n}
\longrightarrow_p0.
\]
The product therefore converges to \(e^{-\kappa d_h(a)}\).
All products are bounded. Joint convergence with an independent test point
implies convergence of \(u_{h,n}\) in \(L^1(P_X)\), in probability over
\(\mathcal D_n\), to one, zero, or \(e^{-\kappa d_h(a)}\) in the respective
regimes.

Finally, conditional on \(\mathcal D_n\), independent tree groups of bounded
size and bounded score predictions give
\[
\left\|\check T_{h,n}-
\E_\Theta(\check T^{\rm tree}_{h,n}\mid\mathcal D_n)\right\|_{L^1(P_X)}
=O_p(B_n^{-1/2})=o_p(1).
\]
Combine this with \eqref{eq:proof-mixture}, the three limits for \(u_{h,n}\),
and the prediction transfer in \eqref{eq:transfers}. This proves
\eqref{eq:selection-limit}, \eqref{eq:recovery-limit}, and
\eqref{eq:intermediate-targets}.

A strictly increasing recoding maps a threshold \(t\) to \(\varphi(t)\)
and preserves \(\varphi(x)\le\varphi(t)\) exactly when \(x\le t\).
All cells and concentration bounds above are expressed in the original
coordinates. Their constants therefore do not depend on the number or
form of the increasing recodings.
\end{proof}

\begin{proof}[Proof of Corollary~\ref{cor:sign}]
In the selection regime, Markov's inequality and
\eqref{eq:selection-limit} give, for \(g=1,2\),
\[
\Prb_X\{\operatorname{sgn}(\widehat T_{g,n})\ne s_g\}
\le\frac{2}{\underline a}
\|\widehat T_{g,n}-a_gs_g\|_{L^1(P_X)}=o_p(1).
\]
The independent signs \(s_1\) and \(s_2\) disagree with probability one half.
On that event, \(|\tau_a|=a_2-a_1\ge\Delta_a\). A second application of
Markov's inequality shows that both fitted magnitudes exceed
\(\underline a/2\) outside a set of \(P_X\)-mass \(o_p(1)\). This proves
\eqref{eq:sign-reversal} and the stated separation from zero.

For any rule \(\pi\),
\[
\operatorname{Reg}_a(\pi)
=\E\!\left[|\tau_a(X)|
\one\{\pi(X)\ne\pi_a^*(X)\}\right].
\]
Since \(a_2>a_1\), \(\pi_a^*=\one\{s_2=1\}\), so the limit of
\(\widehat\pi_{2,n}\) has zero regret. For the other representation,
\[
\operatorname{Reg}_a(\one\{s_1=1\})
=\E\!\left[\tau_a
\{\one\{s_2=1\}-\one\{s_1=1\}\}\right]
=\frac{a_2-a_1}{2}.
\]
Boundedness and classification convergence transfer these values to the
estimated rules, proving \eqref{eq:regret}.
\end{proof}

\begin{proof}[Proof of Corollary~\ref{cor:boundary}]
The capped-Poisson rule satisfies
\[
\E Q_M=\xi_M+e^{-\xi_M}-\E(N_M-M)_+.
\]
For fixed \(\xi_M=\xi\), the final term vanishes as \(M\to\infty\), giving
\(\omega_n\sim(\xi+e^{-\xi})/M_n\). If \(\xi_M\to\infty\) and
\(\xi_M/M\to0\), Poisson concentration gives
\(\E Q_M/\xi_M\to1\). Under the package default,
\(\xi_M/\sqrt M\to1\), so \(\omega_n\sim M_n^{-1/2}\). The two boundaries
follow by substituting these expressions into
\(\mathcal I_n=H_n\omega_n\). For fixed finite \(M\), \(\omega_n>0\) is
constant and \(H_n\to\infty\).
\end{proof}

\begin{proof}[Proof of Proposition~\ref{prop:class}]
Couple the two representations with the common class-indexed randomness in the
proposition. At a common node they draw the same classes, evaluate the same
certified actions, and apply the same scores and tie order. Equal certificates
give identical unordered children, while canonical orientation sends every
point in \(\mathcal X_{\rm cert}\) to the same child. Induction gives identical
tree partitions on the certified array. Leaf observations and forest weights
therefore coincide. Nuisance fits coincide because they are held fixed or
computed from the same canonical matrix with common randomness. The same
inputs, settings, and randomization then give identical stopping decisions,
predictions, variance calculations, rankings, and policy ties.
\end{proof}

\end{appendix}
